\chapter{Methodology}
\label{ch:methodology}

This chapter details each module of RA-GARK in the order they participate in the forward pass: the local view (Section~\ref{sec:local}), the global view's two innovations --- KG-SVD initialization (Section~\ref{sec:kgsvd}) and softmax rationale masking (Section~\ref{sec:softmax}) --- the local-biased fusion gate (Section~\ref{sec:gate}), and the cross-view contrastive regularization (Section~\ref{sec:cl}). Section~\ref{sec:method-train} describes the training procedure and computational complexity. Notation follows Table~\ref{tab:notation}.

\section{Local View: LightGCN}
\label{sec:local}

% TODO: prose for §4.1 Local View
% - K=2 LightGCN propagation, layer-wise average
% - no KG involvement; serves as collaborative-purity anchor
% - cite he2020lightgcn

\section{Global View — KG-SVD Initialization}
\label{sec:kgsvd}

% TODO: prose for §4.2 KG-SVD initialization
% - item--aspect co-occurrence matrix M
% - IDF weighting motivation
% - truncated SVD with k = A * d components
% - reshape to [N_i, A, d]
% - xavier-magnitude rescaling
% - ablation: −5.3% NDCG@20 without SVD init

\begin{figure}[ht]
  \centering
  % \includegraphics[width=0.9\textwidth]{img/kgsvd_init.png}  % TODO: insert figure (b)
  \fbox{\parbox[c][4cm][c]{0.9\textwidth}{\centering\itshape Placeholder: KG-SVD initialization flow\\(item--aspect matrix $\to$ IDF weighting $\to$ truncated SVD $\to$ reshape to $[N_i, A, d]$).}}
  \caption{KG-SVD initialization pipeline. The IDF-weighted item--aspect co-occurrence matrix $\mathbf{M}$ is decomposed by truncated SVD into $A \cdot d$ latent components, which are reshaped into $A$ latent aspect slots of dimension $d$ per item.}
  \label{fig:kgsvd}
\end{figure}

\section{Global View — Softmax Rationale Masking}
\label{sec:softmax}

% TODO: prose for §4.3 softmax rationale masking
% - MLP([u_glo; a_{i,k}]) logit per slot
% - softmax over A slots with temperature τ=0.5
% - weighted sum to form i_glo
% - rationale = which latent slot represents item i for user u
% - design rationale for softmax over sigmoid (competition vs independence)
% - ablation: −7.0% NDCG@20 with sigmoid

\begin{figure}[ht]
  \centering
  % \includegraphics[width=0.9\textwidth]{img/softmax_attention.png}  % TODO: insert figure (c)
  \fbox{\parbox[c][4cm][c]{0.9\textwidth}{\centering\itshape Placeholder: Softmax aspect-saliency attention forward pass\\($u_{\mathrm{glo}}$, $\{a_{i,k}\}_{k=1}^A$ $\to$ MLP logits $\to$ softmax weights $\to$ weighted sum $\to$ $i_{\mathrm{glo}}$).}}
  \caption{Softmax rationale masking. The user-side global embedding $u_{\mathrm{glo}}$ is concatenated with each of the item's $A$ latent aspect slots; a shared MLP produces a logit per slot, which is normalized by softmax with temperature $\tau$. The resulting weights compete over the $A$ slots and produce $i_{\mathrm{glo}}$ as a weighted sum.}
  \label{fig:softmax}
\end{figure}

\section{Local-Biased Fusion Gate}
\label{sec:gate}

% TODO: prose for §4.4 fusion gate
% - independent user_fusion_gate and item_fusion_gate
% - MLP([x_loc; x_glo]) → tanh → linear → sigmoid → α
% - bias init = +5 on final linear; α₀ ≈ σ(5) ≈ 0.993
% - graceful degradation: model starts ≈ pure LightGCN
% - ablation: −5.3% NDCG@20 with bias=0

\begin{figure}[ht]
  \centering
  % \includegraphics[width=0.9\textwidth]{img/fusion_gate.png}  % TODO: insert figure (d)
  \fbox{\parbox[c][4cm][c]{0.9\textwidth}{\centering\itshape Placeholder: Fusion gate signal flow\\($x_{\mathrm{loc}}$, $x_{\mathrm{glo}}$ $\to$ MLP with bias-initialized final linear ($b=+5$) $\to$ sigmoid $\to$ $\alpha$ $\to$ $x_{\mathrm{final}} = \alpha\, x_{\mathrm{loc}} + (1-\alpha)\, x_{\mathrm{glo}}$). Two independent gates: one for $u$, one for $i$.}}
  \caption{Local-biased fusion gate. Independent user- and item-side gates each take the concatenation of their local and global embeddings and produce a scalar $\alpha \in (0, 1)$. The bias of the final linear layer is initialized to $+5$, so that $\alpha^{(0)} \approx 0.993$ and the model begins training as effectively pure LightGCN.}
  \label{fig:gate}
\end{figure}

\section{Cross-View Contrastive Regularization}
\label{sec:cl}

% TODO: prose for §4.5 contrastive learning
% - L_aCL between (i_loc, i_glo); L_uCL between (u_loc, u_glo)
% - InfoNCE form with cl_projector head
% - stop-gradient on global side
% - small fixed weight λ_CL = 0.005 (auxiliary role)

\section{Training and Complexity}
\label{sec:method-train}

% TODO: prose for §4.6 training & complexity
% - Adam, lr, batch size, early stopping
% - complexity breakdown per module (see slide 19 table)
% - inference cost for full-item ranking
